\documentclass[12pt]{article}
\usepackage[utf8]{inputenc}
\usepackage[T1]{fontenc}
\usepackage{hyperref}
\title{Safwa Language Resources for K-12 Curriculum Development}
\author{Compiled by Claude Sonnet 5.5 (claude-sonnet-5-5) for Andr\'as Kornai}
\date{}
\begin{document}
\maketitle

\section{Introduction}
Safwa (Kisafwa, ISO 639-3 sbk, Glottocode safw1238) is a Bantu language of the Mbeya Region of southwest Tanzania, spoken in the Mbozi, Mbeya and Chunya districts. Sources list the dialects Guruka, Mbwila, Poroto and Songwe. Speaker estimates vary widely, from about 100,000 to 322,000. Most speakers also know Swahili. The language is written in the Latin script. Only a small set of resources turned up in search, and none of them is school material.

\section{Primary Educational Resources}
No primer, textbook or school programme for Safwa turned up in search. Omniglot mentions some children's books, but it does not name them and I could not confirm any. See the Omniglot page below.
\begin{itemize}
\item Omniglot, Safwa (Kisafwa): alphabet chart (Herufi za Kisafwa), sample text and background. \url{https://www.omniglot.com/writing/safwa.htm}
\end{itemize}

\section{Linguistic and Reference Materials}
Search did not turn up a dictionary, grammar or formal orthography document. Wikipedia discusses the open question of whether the vowel system has seven vowels or five, which matters for any orthography work.
\begin{itemize}
\item Wikipedia, Safwa language. \url{https://en.wikipedia.org/wiki/Safwa_language}
\item Wikipedia, Safwa people. \url{https://en.wikipedia.org/wiki/Safwa_people}
\item ASJP, Safwa word list. \url{https://asjp.clld.org/languages/SAFWA}
\item Glottolog, Safwa (safw1238). \url{https://glottolog.org/resource/languoid/id/safw1238.bigmap.html}
\end{itemize}

\section{Religious and Community Resources}
A New Testament in Safwa (Biblia Takatifu / Agano Jipya), completed in 2020 by the Door43 Worldwide Mission Community, is on eBible.org under a Creative Commons Attribution 4.0 licence. It is available in several formats, including PDF and ePub. As an openly licensed text it could be reused as a source of connected native-language prose.
\begin{itemize}
\item eBible.org, Kisafwa (sbk) New Testament. \url{https://ebible.org/bible/details.php?id=sbk}
\item eBible.org, Safwa New Testament PDF. \url{https://ebible.org/pdf/sbk/sbk_nt.pdf}
\end{itemize}

\section{Digital and Audio Resources}
Global Recordings Network has a page for Kisafwa (sbk) that lists audio Bible stories and evangelistic programmes. The search summary also mentioned a Jesus Film Project entry, but I did not see that entry myself.
\begin{itemize}
\item Global Recordings Network, Kisafwa. \url{https://globalrecordings.net/en/language/sbk}
\end{itemize}

\section{Curriculum Development Context}
The reference description says Safwa is not known to be taught in schools, and nothing found contradicts that. A curriculum would have to start almost from nothing. The Latin-script alphabet chart, the openly licensed New Testament and the audio recordings are the only raw materials found, and the vowel-system question would need settling before an orthography is standardised. Swahili is the national school language, so a bilingual transitional model is the likely route. The dossier is thin because little exists in what was searched. Local sources in Tanzania may hold more.
\end{document}